\documentclass[11pt]{article}

\usepackage[margin=1in]{geometry}
\usepackage{booktabs}
\usepackage{graphicx}
\usepackage[hidelinks]{hyperref}
\usepackage{microtype}
\usepackage{amsmath}
\usepackage{enumitem}
\usepackage[table]{xcolor}
\usepackage{float}
\usepackage{tabularx}
\usepackage{url}

\title{Rewrite or Patch? A Controlled Study of Identity-Preserving State Mutation Interfaces}
\author{Yin Li\thanks{Corresponding author. Email: \texttt{kxl474@student.bham.ac.uk}}\\
University of Birmingham, Edgbaston, Birmingham B15 2TT, United Kingdom\\
\texttt{kxl474@student.bham.ac.uk}}
\date{August 2026}

\newcommand{\orderdelta}{\textsc{OrderDeltaBench}}
\emergencystretch=2em

\begin{document}
\maketitle

\begin{abstract}
Structured-output interfaces make language-model tool calls parseable, but
transaction agents must also update persistent state correctly.
A restaurant ordering agent rarely receives only one clean order: users add
items, remove one of several similar lines, change modifiers, reverse earlier
decisions, and introduce allergy or dietary constraints after a cart already
exists. We introduce \orderdelta, a deterministic benchmark for stateful cart
edits in which each model must transform a current order into a safe updated
transaction. The archived v2 experiment compares two executable interface
bundles: an ID-free full-order rewrite and stable-line operations. The promoted
v3 experiment removes that identity-information confound and adds an
identity-aware verifier, authorized-write diffs, atomic execution, mutation
testing, and three repeated provider requests. Across 22,680 recorded calls to
seven models and three strict structured-output interfaces, typed patches have
a model-dependent effect relative to ID-bearing rewrites. Qwen3-32B improves by
20.3 percentage points (95\% cluster CI [12.4, 28.4]) and
Cosmos3-Super-Reasoner by 12.7 points ([6.7, 19.4]); both remain significant
after Benjamini--Hochberg adjustment across seven model-specific primary tests.
The other five adjusted comparisons are not significant, and Hermes-4-70B has
a negative point estimate. A generic JSON Patch contract records no complete
semantic successes because every schema-valid response misses the required
version precondition; this is a contract-specific negative result. The final
219 GLM-5.1 calls fail with HTTP 402 after the supplied provider budget is
exhausted and are retained as balanced unconditional failures. These results
support an interface-by-model interaction, not a universal patch advantage.
\end{abstract}

\section*{Current Validity Notice}
The controlled v3 result is the primary evidence. The later ``Archived v2
Results'' section is retained to document the original 5,040-call artifact but
must not be read causally: v2 removes stable line IDs only from rewrite prompts,
and its multiset verifier can miss wrong-object writes when values collide. V3
uses information-equivalent \texttt{rewrite\_with\_ids}, typed patch, and JSON
Patch conditions; atomic execution; authorized JSON-pointer diffs; mutation
testing; clustered inference; and explicit repeated-request reliability.

\noindent\textbf{Keywords:} large language models; tool use; structured output;
stateful agents; transaction safety; JSON Patch; empirical software engineering

\section{Introduction}
LLM agents are increasingly used as software components that translate natural
language into executable API calls. In transactional systems, the risky output
is often not malformed text. It is a well-formed object that changes the wrong
state. A user may say ``remove the plain vegan bowl'', ``put the cheese back
but keep no onion'', or ``I forgot to say I have a peanut allergy'' after the
cart already contains multiple lines. A safe system must identify the intended
line, preserve unrelated lines, update constraints, and refuse or clarify when
the requested edit would make the transaction unsafe.

This paper studies that state-update layer. Prior work on structured outputs
and function calling has improved syntactic adherence to schemas
\cite{openai2024structured,jsonschema2020,patil2025bfcl,yang2025structeval}.
Agent benchmarks have also moved toward interactive, stateful tasks
\cite{yao2024taubench,lu2024toolsandbox,enterpriseopsgym2026}. However, a
deployment engineer still faces a narrower interface-design question: when an
LLM updates a persistent transaction, should it regenerate the whole resource
or return a minimal patch over stable object identifiers?

\orderdelta{} isolates this question in restaurant ordering. The domain is
small enough to verify exactly, but rich enough to contain common production
failure modes: duplicate lines, scoped modifiers, quantity deltas, reversal of
previous modifiers, unavailable catalog entries, stale references to absent
cart lines, and user constraints that newly conflict with the current cart.
The benchmark compares:
\begin{itemize}[leftmargin=*]
  \item \textbf{ID-bearing full-order rewrite}: the model receives stable line
  IDs and returns the entire updated order while preserving identity.
  \item \textbf{Typed line patch}: the model receives the same state and
  returns typed operations such as \texttt{add\_item}, \texttt{remove\_line},
  \texttt{update\_line}, and \texttt{set\_constraints}.
  \item \textbf{JSON Patch}: the model returns an RFC 6902-style operation list
  with a required optimistic-concurrency version test.
\end{itemize}

The study contributes:
\begin{itemize}[leftmargin=*]
  \item a deterministic 360-case benchmark for stateful transactional edits;
  \item three information-equivalent strict JSON Schema interface contracts;
  \item an identity-aware executable verifier with atomic application,
  authorized-write diffs, mutation testing, and fail-closed safety checks;
  \item a controlled comparison across seven models and 22,680 recorded calls,
  with cluster-aware inference and repeated-request reliability.
\end{itemize}

\section{Related Work}
\paragraph{Structured output and function calling.}
JSON Schema and provider structured-output modes define the structural contract
between model and software system \cite{jsonschema2020,openai2024structured}.
BFCL evaluates whether models can select functions and produce suitable
arguments \cite{patil2025bfcl}, while Gorilla exposed hallucinated APIs as a
core limitation of tool use \cite{patil2023gorilla}. StructEval broadens the
view of structured generation across renderable and non-renderable formats
\cite{yang2025structeval}. \orderdelta{} assumes strict schemas are useful, but
measures the next layer: whether the resulting object correctly updates a
transactional state.

\paragraph{Stateful tool-use benchmarks.}
\(\tau\)-bench and ToolSandbox evaluate agents in interactive domains with
state and user interaction \cite{yao2024taubench,lu2024toolsandbox}. WebMall and
ShoppingComp focus on shopping-agent retrieval, comparison, and decision
quality \cite{webmall2026,tou2025shoppingcomp}. EnterpriseOps-Gym evaluates
long-horizon enterprise tool use with persistent databases and policy
constraints \cite{enterpriseopsgym2026}. SWE-bench-Live illustrates the same
methodological direction in software engineering: scalable, executable,
continuously refreshed evaluation with explicit environment validation
\cite{zhang2025swebenchlive}. These benchmarks are broader than
\orderdelta. Our benchmark is deliberately surgical: it fixes the domain and
prompt context so that interface choice and state-update semantics can be
measured case by case.

\paragraph{Patch-based editing.}
JSON Patch defines a sequence of operations for partial updates to JSON
documents \cite{rfc6902}. Recent LLM work on JSON editing argues that patch
generation can reduce token cost, and that stable keys help avoid array-index
errors \cite{duanis2025jsonwhisperer}. \orderdelta{} adapts that intuition to
transactional agents: the key question is not only efficient editing, but
whether line-keyed operations reduce wrong-line edits, cart drift, and unsafe
state changes.

\paragraph{Reporting empirical LLM studies.}
Because model APIs, model snapshots, and provider behavior change over time,
empirical LLM studies need explicit model identifiers, prompts, logs,
configuration, metrics, and limitations \cite{baltes2025llmguidelines}. The
artifact records exact provider model IDs, raw responses, usage fields,
latency, prompts, generated cases, and verifier outputs.

\section{Benchmark Design}
\subsection{Domain Contract}
\orderdelta{} reuses a compact restaurant menu with seven SKUs: classic burger,
vegan bowl, chicken wrap, satay noodles, margherita pizza, fries, and cola.
Each SKU has allowed sizes, default modifiers, removable modifiers, addable
modifiers, allergens, and dietary tags. A cart line contains a stable
\texttt{line\_id}, SKU, quantity, size, additions, removals, and free-form
special instructions. A cart also contains stated allergen and dietary
constraints.

Each benchmark case has a current order, a short edit history, a user edit
utterance, an expected status, an expected final order, and audit metadata
marking target and unchanged line IDs. The edit history makes the prompt more
like a stateful transaction, while the current order remains the authoritative
state for scoring. The expected status is \texttt{accepted},
\texttt{needs\_clarification}, or \texttt{rejected\_safety}. When an edit is
not accepted, the verifier expects fail-closed behavior: the model should not
execute the requested cart edit.

\subsection{Categories}
The 360 cases are evenly distributed across 12 categories:
\begin{enumerate}[leftmargin=*]
  \item adding a valid item to an existing cart;
  \item removing one line among similar duplicate lines;
  \item increasing a line quantity;
  \item decreasing a line quantity;
  \item adding a modifier to only one referenced line;
  \item reversing one previous modifier decision while keeping another;
  \item changing the size of one duplicate line;
  \item replacing an item while preserving sides or drinks;
  \item adding a safe allergen or dietary constraint;
  \item introducing a constraint that conflicts with the current cart;
  \item requesting an unavailable item or modifier;
  \item referring to a stale line that is not in the current cart.
\end{enumerate}
Each category contains 30 cases generated deterministically from eight
hand-written semantic templates plus lexical prefix and suffix variation. No
model is used to create oracle labels.

\subsection{Interfaces}
The full-order rewrite interface asks the model to return:
\begin{center}
\small
\texttt{\{status, updated\_order, clarification\_question, reasons\}}
\end{center}
The current cart is shown with item positions for human-readable reference, but
the output schema forbids positions and line IDs.

The line-ID patch interface asks the model to return:
\begin{center}
\small
\texttt{\{status, operations, clarification\_question, reasons\}}
\end{center}
Each operation has an operation type, optional line ID, optional item, optional
quantity and size updates, modifier add/remove lists, modifier reversal fields,
and optional constraints. This is not a direct RFC 6902 schema; it is a
domain-specific patch contract inspired by JSON Patch but designed to be
auditable and executable against cart lines.

\section{Evaluation Method}
For every model, interface, and case, the runner sends one Nebius
OpenAI-compatible chat-completion request at temperature zero with strict
\texttt{json\_schema} response formatting \cite{nebius2026studio}. The
evaluated models are Nebius-accessible open-weight models rather than a general
leaderboard over all commercial systems. The raw result records provider
success, response content, finish reason, usage, and latency.

The verifier first checks JSON extraction and interface schema validity. For
rewrite outputs, it directly scores the returned \texttt{updated\_order}. For
patch outputs, it applies the operation list to the current cart, then strips
line IDs and compares the resulting order with the oracle final order. The
main metrics are:
\begin{itemize}[leftmargin=*]
  \item \textbf{Semantic success}: schema-valid output, correct status, exact
  final item multiset, exact constraints, executable patch operations, and
  fail-closed state preservation for non-accepted edits.
  \item \textbf{Unsafe state change}: accepting or changing state when the
  oracle requires clarification or safety rejection, or accepting a final order
  that conflicts with stated constraints.
  \item \textbf{Unintended drift}: changing a line that the oracle marks as
  unchanged.
  \item \textbf{Identity error}: a wrong-line error in a case with a target
  line.
\end{itemize}

\subsection{Scoring Fairness}
Both interfaces are reduced to the same object before semantic scoring: a
status decision and a final cart state. For rewrite outputs, the final cart is
the returned \texttt{updated\_order}. For patch outputs, the verifier applies
the returned operations to the current cart, strips line identifiers, and then
compares the resulting cart to the same oracle used for rewrite scoring.

The verifier performs limited JSON extraction only: it accepts a raw JSON
object, a fenced JSON object, or the substring from the first opening brace to
the last closing brace. It does not repair missing fields, bad operation names,
wrong primitive types, invented modifiers, or malformed item objects.
Schema-invalid outputs receive \texttt{semantic\_ok=false}, while component
metrics still record any parsable partial behavior. Hallucinated cart lines,
patches that reference absent line IDs, or stale-reference edits that mutate
the cart are treated as execution failures. For non-accepted oracle cases, a
model must fail closed: the status must match the oracle and the requested cart
edit must not be executed. This normalization does not remove the
input-information confound: the patch condition exposes stable IDs while the
v2 rewrite condition does not. Therefore the reported paired comparison
estimates the effect of the full interface bundle, not patch syntax alone.

For paired comparisons between rewrite and line-patch interfaces, the same
case is evaluated under both interfaces for the same model. We report the
patch-minus-rewrite rate difference, paired bootstrap 95\% confidence intervals
over cases, and an exact McNemar test on discordant paired outcomes.

\section{Controlled V3 Experiment}
V3 evaluates 360 cases under \texttt{rewrite\_with\_ids}, a tagged-union typed
patch, and JSON Patch. Seven catalog models receive all three identity-bearing
conditions in counterbalanced order at temperature zero and a 2,400-token
output limit. Each model/interface/case cell is requested three times. The
three calls measure serving and output reliability; they are not treated as
independent task designs.

All 22,680 expected keys are recorded with zero duplicate keys or protocol
mismatches. Independent evaluator replay reproduces all 22,680 stored scores.
Six models complete all calls. The final 219 GLM-5.1 calls fail after all
configured retries with HTTP 402 because the supplied provider budget is
exhausted. These failures are balanced across interfaces (73 each), remain
unconditional failures, and are also analyzed in a provider-success-only
sensitivity analysis. Returned token usage and the frozen catalog prices imply
an estimated cost of \$24.3292; this value is not a billing invoice.

The primary estimand is the typed-patch-minus-ID-rewrite difference in
unconditional semantic success. Outcomes are averaged within each of 96
hand-written \texttt{semantic\_program\_id} clusters. We report a cluster
bootstrap interval, a cluster sign-flip test, and Benjamini--Hochberg
adjustment across the seven model-specific primary comparisons. Case-level
McNemar values are retained in the artifact as descriptive diagnostics only.

\begin{table}[H]
\centering
\scriptsize
\resizebox{\linewidth}{!}{\begin{tabular}{llrrrrr}
\toprule
Model & Metric & Rewrite+IDs & Typed patch & $\Delta$ [95\% cluster CI] & $p$ & BH $q$ \\
\midrule
Qwen3-235B-A22B & Semantic & 80.2 & 85.9 & +5.4 [-0.4, +11.8] & 0.089 & 0.125 \\
Qwen3-235B-A22B & Unsafe & 2.6 & 1.3 & -1.2 [-3.3, +0.3] & 0.375 & 1.000 \\
Qwen3-235B-A22B & Drift & 0.1 & 1.4 & +1.3 [+0.1, +3.0] & 0.125 & 0.219 \\
Qwen3-32B & Semantic & 65.5 & 85.3 & +20.3 [+12.4, +28.4] & <.001 & <.001 \\
Qwen3-32B & Unsafe & 14.5 & 5.4 & -9.5 [-14.9, -4.7] & <.001 & 0.002 \\
Qwen3-32B & Drift & 12.4 & 3.0 & -10.0 [-15.8, -4.9] & <.001 & 0.002 \\
gpt-oss-120b & Semantic & 91.3 & 91.9 & +0.1 [-4.0, +4.4] & 0.969 & 0.969 \\
gpt-oss-120b & Unsafe & 1.2 & 1.6 & +0.6 [-1.3, +2.8] & 0.625 & 1.000 \\
gpt-oss-120b & Drift & 0.6 & 1.6 & +1.1 [-0.1, +3.0] & 0.375 & 0.438 \\
Hermes-4-70B & Semantic & 74.1 & 68.6 & -5.4 [-13.5, +2.6] & 0.198 & 0.231 \\
Hermes-4-70B & Unsafe & 1.1 & 1.2 & +0.1 [-2.6, +2.2] & 1.000 & 1.000 \\
Hermes-4-70B & Drift & 0.0 & 1.2 & +1.2 [+0.1, +2.6] & 0.125 & 0.219 \\
Hermes-4-405B & Semantic & 82.6 & 88.6 & +6.2 [-0.5, +13.0] & 0.070 & 0.122 \\
Hermes-4-405B & Unsafe & 1.9 & 1.9 & +0.0 [+0.0, +0.0] & 1.000 & 1.000 \\
Hermes-4-405B & Drift & 0.0 & 2.2 & +2.2 [+0.3, +4.8] & 0.062 & 0.219 \\
Cosmos3-Super-Reasoner & Semantic & 69.3 & 81.9 & +12.7 [+6.7, +19.4] & <.001 & <.001 \\
Cosmos3-Super-Reasoner & Unsafe & 1.2 & 0.3 & -1.0 [-3.1, +0.5] & 0.500 & 1.000 \\
Cosmos3-Super-Reasoner & Drift & 1.4 & 0.2 & -1.4 [-3.5, +0.2] & 0.250 & 0.350 \\
GLM-5.1 & Semantic & 79.4 & 85.7 & +6.2 [+0.7, +12.6] & 0.042 & 0.097 \\
GLM-5.1 & Unsafe & 0.1 & 0.1 & +0.0 [-0.3, +0.3] & 1.000 & 1.000 \\
GLM-5.1 & Drift & 0.1 & 0.1 & +0.0 [-0.3, +0.3] & 1.000 & 1.000 \\
\bottomrule
\end{tabular}
}
\caption{Controlled v3 comparisons. Rates and effects are percentage points.
Inference clusters at the semantic-program level; BH adjustment is within each
metric and comparison family.}
\label{tab:v3-controlled}
\end{table}

Qwen3-32B improves by 20.3 points (95\% cluster CI [12.4, 28.4], adjusted
\(q=.00035\)) and Cosmos3-Super-Reasoner by 12.7 points ([6.7, 19.4],
\(q=.00035\)). The adjusted comparisons for Qwen3-235B-A22B,
gpt-oss-120b, both Hermes deployments, and GLM-5.1 are not significant.
GLM-5.1's provider-success-only effect is +5.9 points ([0.3, 12.4]), with
cluster \(p=.0564\) and adjusted \(q=.1219\). The evidence therefore supports
a model-dependent interface effect rather than a general patch advantage.

Repeated requests also expose output variability hidden by one-shot evaluation.
For typed patch, all-three semantic success ranges from 65.3\% for Hermes-4-70B
to 89.7\% for gpt-oss-120b among the six fully available models. Byte-identical
three-run output can be much lower: 2.2\% for gpt-oss-120b and 6.4\% for
Qwen3-235B-A22B typed patches, even though semantic-outcome agreement remains
at least 93.1\%. GLM reliability is separately confounded by the balanced
provider-budget failures.

The generic JSON Patch condition has zero semantic successes for every model.
All 7,353 schema-valid JSON Patch responses omit or miss the required first
\texttt{test /version} operation; remaining responses are truncated, malformed,
empty, or provider failures. This observation applies only to the exact schema,
prompt rules, deployments, and output budget recorded here.

\section{Archived V2 Results}
Table~\ref{tab:main} gives the main results. All 5,040 provider calls
completed. All evaluated models except Gemma-2-2B produced near-universal
schema-valid outputs in both interfaces; Gemma-2-2B is retained as a low-end
counterexample because line patch raised schema validity from 3.6\% to 70.0\%
without producing exact semantic success.

\begin{table}[H]
\centering
\small
\resizebox{\linewidth}{!}{\begin{tabular}{llrrrrrrr}
\toprule
Model & Interface & $n$ & Provider & Schema & Semantic & Unsafe & Drift & Identity \\
\midrule
Qwen3-235B-A22B & rewrite & 360 & 100.0 & 100.0 & 86.1 & 1.9 & 0.0 & 0.0 \\
Qwen3-235B-A22B & line patch & 360 & 100.0 & 100.0 & 83.9 & 0.6 & 0.8 & 0.8 \\
Qwen3-32B & rewrite & 360 & 100.0 & 100.0 & 64.7 & 21.9 & 5.8 & 0.8 \\
Qwen3-32B & line patch & 360 & 100.0 & 100.0 & 80.8 & 7.2 & 4.4 & 0.0 \\
Qwen3-30B-A3B & rewrite & 360 & 100.0 & 100.0 & 74.2 & 6.9 & 3.6 & 0.0 \\
Qwen3-30B-A3B & line patch & 360 & 100.0 & 100.0 & 70.0 & 7.2 & 5.3 & 0.8 \\
Llama-3.3-70B & rewrite & 360 & 100.0 & 100.0 & 76.7 & 3.1 & 4.4 & 1.1 \\
Llama-3.3-70B & line patch & 360 & 100.0 & 100.0 & 72.8 & 2.8 & 2.8 & 0.0 \\
Llama-3.1-8B & rewrite & 360 & 100.0 & 100.0 & 25.3 & 9.7 & 29.4 & 15.8 \\
Llama-3.1-8B & line patch & 360 & 100.0 & 100.0 & 48.1 & 6.9 & 11.4 & 3.6 \\
Gemma-2-2B & rewrite & 360 & 100.0 & 3.6 & 0.0 & 1.9 & 0.8 & 0.0 \\
Gemma-2-2B & line patch & 360 & 100.0 & 70.0 & 0.0 & 11.1 & 0.0 & 0.0 \\
Nemotron-Ultra-253B & rewrite & 360 & 100.0 & 100.0 & 79.7 & 0.0 & 0.0 & 0.0 \\
Nemotron-Ultra-253B & line patch & 360 & 100.0 & 99.7 & 78.3 & 0.0 & 0.3 & 0.3 \\
\bottomrule
\end{tabular}
}
\caption{Main results. Values other than \(n\) are percentages. ``Semantic''
is exact state-update success; ``Unsafe'' is an unsafe accepted or changed
state; ``Drift'' is unintended change to protected cart lines.}
\label{tab:main}
\end{table}

\subsection{Line IDs Help Where Identity Is the Bottleneck}
The clearest positive effects appear for Qwen3-32B and Llama-3.1-8B. For
Qwen3-32B, line patches improve semantic success from 64.7\% to 80.8\%, a
+16.1 point paired difference with a 95\% confidence interval of [+11.1,
+21.1] and McNemar \(p<.001\). The same interface reduces unsafe state changes
from 21.9\% to 7.2\%. For Llama-3.1-8B, patching improves semantic success
from 25.3\% to 48.1\%, a +22.8 point difference with a 95\% confidence
interval of [+16.9, +28.9] and McNemar \(p<.001\). It also cuts unintended
drift from 29.4\% to 11.4\%. These are the strongest evidence that stable line
identifiers help when a model would otherwise rewrite unrelated cart state.

Table~\ref{tab:paired} shows that the gain is not universal. Qwen3-235B-A22B
slightly favors rewrite on exact semantic success (86.1\% versus 83.9\%),
Qwen3-30B-A3B also favors rewrite (74.2\% versus 70.0\%), Llama-3.3-70B favors
rewrite (76.7\% versus 72.8\%), and Nemotron-Ultra-253B is nearly tied
(79.7\% versus 78.3\%). These differences are small relative to the large
positive effects for Qwen3-32B and Llama-3.1-8B.

\begin{table}[H]
\centering
\small
\resizebox{\linewidth}{!}{\begin{tabular}{llrrrr}
\toprule
Model & Metric & Rewrite & Patch & $\Delta$ [95\% CI] & $p$ \\
\midrule
Qwen3-235B-A22B & Semantic & 86.1 & 83.9 & -2.2 [-4.7, +0.3] & 0.134 \\
Qwen3-235B-A22B & Unsafe & 1.9 & 0.6 & -1.4 [-2.8, +0.0] & 0.125 \\
Qwen3-235B-A22B & Drift & 0.0 & 0.8 & +0.8 [+0.0, +1.9] & 0.250 \\
Qwen3-32B & Semantic & 64.7 & 80.8 & +16.1 [+11.1, +21.1] & <.001 \\
Qwen3-32B & Unsafe & 21.9 & 7.2 & -14.7 [-18.6, -10.8] & <.001 \\
Qwen3-32B & Drift & 5.8 & 4.4 & -1.4 [-3.3, +0.6] & 0.267 \\
Qwen3-30B-A3B & Semantic & 74.2 & 70.0 & -4.2 [-8.6, +0.3] & 0.091 \\
Qwen3-30B-A3B & Unsafe & 6.9 & 7.2 & +0.3 [-1.1, +1.7] & 1.000 \\
Qwen3-30B-A3B & Drift & 3.6 & 5.3 & +1.7 [+0.0, +3.3] & 0.109 \\
Llama-3.3-70B & Semantic & 76.7 & 72.8 & -3.9 [-7.8, +0.0] & 0.076 \\
Llama-3.3-70B & Unsafe & 3.1 & 2.8 & -0.3 [-0.8, +0.0] & 1.000 \\
Llama-3.3-70B & Drift & 4.4 & 2.8 & -1.7 [-3.1, -0.6] & 0.031 \\
Llama-3.1-8B & Semantic & 25.3 & 48.1 & +22.8 [+16.9, +28.9] & <.001 \\
Llama-3.1-8B & Unsafe & 9.7 & 6.9 & -2.8 [-5.6, +0.0] & 0.076 \\
Llama-3.1-8B & Drift & 29.4 & 11.4 & -18.1 [-23.1, -13.3] & <.001 \\
Gemma-2-2B & Semantic & 0.0 & 0.0 & +0.0 [+0.0, +0.0] & 1.000 \\
Gemma-2-2B & Unsafe & 1.9 & 11.1 & +9.2 [+6.1, +12.2] & <.001 \\
Gemma-2-2B & Drift & 0.8 & 0.0 & -0.8 [-1.9, +0.0] & 0.250 \\
Nemotron-Ultra-253B & Semantic & 79.7 & 78.3 & -1.4 [-4.2, +1.4] & 0.458 \\
Nemotron-Ultra-253B & Unsafe & 0.0 & 0.0 & +0.0 [+0.0, +0.0] & 1.000 \\
Nemotron-Ultra-253B & Drift & 0.0 & 0.3 & +0.3 [+0.0, +0.8] & 1.000 \\
\bottomrule
\end{tabular}
}
\caption{Paired rewrite-versus-patch comparisons. \(\Delta\) is line patch
minus rewrite in percentage points.}
\label{tab:paired}
\end{table}

\subsection{Patch Interfaces Are Not Universally Better}
Category-level results in Table~\ref{tab:category} explain the mixed aggregate
effect. Line patches help on duplicate-line removal (51.0\% to 81.0\%),
quantity increases (73.8\% to 81.9\%), scoped add-ons (69.0\% to 80.0\%), and
scoped size edits (64.3\% to 74.3\%). These are exactly the categories where
stable line identity should reduce accidental whole-cart rewriting.

However, patching hurts modifier reversal (71.9\% to 49.5\%) and slightly
hurts item replacement (51.4\% to 48.1\%). In modifier-reversal cases,
the model must distinguish adding a modifier from removing an item from a
previous \texttt{remove} list, or distinguish removing a previous add-on from
adding a default removal. The patch interface makes the operation explicit, but
also introduces operation-selection complexity. This is an important negative
result: patch interfaces are a tool for identity and locality, not a general
semantic reasoning solution.

\begin{table}[H]
\centering
\small
\begin{tabular}{llrrr}
\toprule
Category & Interface & $n$ & Semantic & Unsafe \\
\midrule
add item & rewrite & 210 & 66.7 & 5.7 \\
add item & line patch & 210 & 74.8 & 0.0 \\
remove duplicate & rewrite & 210 & 51.0 & 0.0 \\
remove duplicate & line patch & 210 & 81.0 & 0.0 \\
qty increase & rewrite & 210 & 73.8 & 0.5 \\
qty increase & line patch & 210 & 81.9 & 0.0 \\
qty decrease & rewrite & 210 & 77.1 & 0.0 \\
qty decrease & line patch & 210 & 80.0 & 0.0 \\
scoped add-on & rewrite & 210 & 69.0 & 4.3 \\
scoped add-on & line patch & 210 & 80.0 & 0.0 \\
modifier reversal & rewrite & 210 & 71.9 & 0.5 \\
modifier reversal & line patch & 210 & 49.5 & 0.0 \\
scoped size & rewrite & 210 & 64.3 & 2.9 \\
scoped size & line patch & 210 & 74.3 & 0.0 \\
replace item & rewrite & 210 & 51.4 & 11.0 \\
replace item & line patch & 210 & 48.1 & 0.0 \\
safe constraint & rewrite & 210 & 36.2 & 0.0 \\
safe constraint & line patch & 210 & 39.5 & 0.0 \\
new conflict & rewrite & 210 & 9.0 & 11.0 \\
new conflict & line patch & 210 & 0.0 & 14.8 \\
unavailable edit & rewrite & 210 & 51.0 & 32.9 \\
unavailable edit & line patch & 210 & 58.1 & 31.4 \\
stale reference & rewrite & 210 & 75.7 & 9.5 \\
stale reference & line patch & 210 & 76.7 & 15.2 \\
\bottomrule
\end{tabular}

\caption{Semantic success and unsafe state-change rate by category, aggregated
across the seven evaluated models.}
\label{tab:category}
\end{table}

\subsection{Safety and Catalog Decisions Remain Hard}
The hardest categories are safety and catalog boundary cases. Newly conflicting
constraints have 9.0\% semantic success under rewrite and 0.0\% under line
patch. Unavailable edits have high unsafe state-change rates: 32.9\% under
rewrite and 31.4\% under line patch. Stale references also remain risky:
semantic success is similar under rewrite and line patch (75.7\% versus
76.7\%), but unsafe state changes are higher under line patch (15.2\% versus
9.5\%).

This result matters because these failures occur after structured-output
control. A schema can force the model to choose a valid status string and emit
well-typed operations, but it cannot decide whether a newly stated peanut
allergy invalidates existing satay noodles, whether mozzarella is unavailable
for a vegan bowl, or whether the user is referring to a line that no longer
exists.

\subsection{Error Taxonomy}
Table~\ref{tab:taxonomy} separates the main failure types. Aggregated across
models, line patch reduces unsafe state changes from 6.5\% to 5.1\%, drift
from 6.3\% to 3.6\%, identity errors from 2.5\% to 0.8\%, status errors from
27.3\% to 17.7\%, order errors from 28.9\% to 24.0\%, and constraint errors
from 25.7\% to 16.7\%. Operation errors remain a patch-specific failure mode,
but they are lower than rewrite schema and extraction failures because most
patch-capable models produce executable operations.

\begin{table}[H]
\centering
\small
\resizebox{\linewidth}{!}{\begin{tabular}{lrrrrrrrr}
\toprule
Interface & $n$ & Unsafe & Drift & Identity & Status & Order & Constraint & Operation \\
\midrule
line patch & 2520 & 5.1 & 3.6 & 0.8 & 17.7 & 24.0 & 16.7 & 4.3 \\
rewrite & 2520 & 6.5 & 6.3 & 2.5 & 27.3 & 28.9 & 25.7 & 13.8 \\
\bottomrule
\end{tabular}
}
\caption{Multi-label error taxonomy by interface. Values are
percentages of cases. Labels do not sum to 100\%.}
\label{tab:taxonomy}
\end{table}

\section{Discussion}
The controlled result narrows the engineering recommendation. Typed patches
materially improve two deployments after cluster-level multiplicity control,
but the effect is near zero, uncertain, or negative elsewhere. Interface
locality is therefore a model-dependent reliability control, not a substitute
for model evaluation or deterministic transaction verification.

\subsection{Interface Design Is a Reliability Control}
The Qwen3-32B and Llama-3.1-8B results suggest a concrete engineering
recommendation: if an LLM must edit persistent transactional state, the API
should expose stable line IDs and prefer localized operations for
identity-sensitive edits. This reduces the surface area over which the model
can accidentally rewrite unrelated state. The result is consistent with the
general motivation for patch formats and stable keys
\cite{rfc6902,duanis2025jsonwhisperer}, but the evidence here is transactional
rather than purely document-editing oriented.

\subsection{Patch APIs Need Domain Verifiers}
The same results warn against treating patching as a complete guardrail.
Patches improve locality, but unavailable edits and newly conflicting
constraints are still semantic policy decisions. In a production ordering
system, the LLM output should be treated as a candidate state transition. A
deterministic verifier should check catalog membership, modifier availability,
line existence, allergen and dietary constraints, and whether non-accepted
edits leave state unchanged.

\subsection{A Strong Model Narrows the Interface Gap}
Qwen3-235B-A22B, Llama-3.3-70B, and Nemotron-Ultra-253B perform similarly
under both interfaces. This does not make interface design irrelevant. It means
that when the model is strong enough for a given edit distribution, line IDs
mainly preserve good behavior rather than create large gains. For mid-sized or
weaker models, the interface can materially change the error profile.

\section{Threats to Validity}
\paragraph{Synthetic domain.}
The benchmark uses a synthetic restaurant menu. This improves auditability and
exact verification, but real ordering systems have larger catalogs, prices,
inventory, tax, substitutions, user profiles, and payment state.

\paragraph{Single-edit setting.}
Each case evaluates one edit against a current cart with a short edit history
rather than a fully interactive long conversation. This isolates state-update
semantics, but does not measure dialogue recovery, memory across many turns, or
user clarification behavior.

\paragraph{Benchmark size and templating.}
The v2 and v3 datasets have 360 deterministic cases. This is larger and more varied
than the initial 120-case artifact, but it is still template-generated over a
compact synthetic menu. A stronger version should add multiple menus, more
catalog variation, paraphrases from real ordering logs where available, and
interactive clarification turns.

\paragraph{Oracle strictness.}
V2 semantic success requires exact equality over item multisets and constraints.
Some deviations may be recoverable in a real interface, but exactness is
appropriate for a transaction contract that may be sent directly to an API.
However, item multisets cannot establish identity preservation when two lines
have equal values or converge to equal values. V3 replaces this check with
stable-ID state diffs and reports identity only for identity-observable modes.

\paragraph{Statistical dependence.}
Lexical variants derived from one semantic template are not independent task
designs. The archived tables use case-level paired bootstrap intervals; v3 also
reports bootstrap intervals and sign-flip tests over 96
\texttt{semantic\_program\_id} clusters, with Benjamini--Hochberg adjustment
across model-specific primary comparisons. Three repeated provider requests
measure reliability and do not triple the number of independent task designs.

\paragraph{Model and provider instability.}
The v3 model IDs were available through Nebius at run time on August 31, 2026.
Provider routing, model snapshots, response-format support, and model behavior
can change. Three deployments return no model fingerprint. The final 219
GLM-5.1 calls fail after the supplied budget is exhausted; balanced
unconditional and provider-success sensitivity analyses are both required.
The raw logs record exact model IDs, request IDs, timestamps, outputs, and the
available fingerprints to support replication.

\paragraph{Prompt and schema choices.}
The study compares three concrete interface designs, not every possible rewrite
or patch prompt. In particular, the generic JSON Patch failure is inseparable
from its exact schema and prompt rules. A separately labeled contract ablation
is needed before drawing broader conclusions about RFC 6902.

\section{Conclusion}
\orderdelta{} shows that stateful transactional reliability is not solved by
valid JSON. After equalizing identity information, typed patches produce large,
cluster-significant gains for Qwen3-32B and Cosmos3-Super-Reasoner, but not for
the other five models after multiplicity adjustment. The result is an
interface-by-model interaction, not a universal patch advantage. Repeated calls
also show that temperature-zero output is often not byte deterministic. A safe
transaction system should therefore choose interfaces empirically, retain
stable object identity, and enforce atomic domain verification before any
model-produced state transition is executed.

\section*{Data and Code Availability}
The dataset, prompts, runner, verifier, raw model outputs, analysis scripts,
and generated tables are archived on Zenodo
\cite{li2026orderdeltaartifact} at
\url{https://doi.org/10.5281/zenodo.20236620}. The development repository is
available at
\url{https://github.com/yinli-systems/order-delta-bench}. The controlled raw
records are in \path{results/provider_v3/formal/raw/*.jsonl.gz}, with structural,
source, evaluator-replay, reliability, cost, and sensitivity receipts in the
same directory. The archived Zenodo release predates this v3 provider run. The
repository does not store provider API keys; keys are read only from the
execution environment.

\section*{Funding}
This research did not receive any specific grant from funding agencies in the
public, commercial, or not-for-profit sectors.

\section*{Declaration of Competing Interest}
The author declares no known competing financial interests or personal
relationships that could have appeared to influence the work reported in this
paper.

\section*{CRediT Author Statement}
Yin Li: Conceptualization, Methodology, Software, Validation, Formal analysis,
Investigation, Data curation, Writing -- original draft, Writing -- review and
editing, Visualization.

\section*{Declaration of Generative AI and AI-Assisted Technologies in the Manuscript Preparation Process}
During the preparation of this work the author used OpenAI Codex/GPT-5 to
assist with software implementation, experiment orchestration, literature
organization, and manuscript drafting. The author reviewed and edited the
outputs and takes full responsibility for the content of the manuscript and
artifact.

\section*{Acknowledgements}
No additional acknowledgements.

\bibliographystyle{plain}
\bibliography{references}

\end{document}
